\begin{theorem}[\texorpdfstring{$2\pi i\ZZ$}{2pi i Z}-晶格分支点与非 \texorpdfstring{$D$}{D}-有限性]\label{thm:fold-gauge-anomaly-singularity-lattice-nonholonomic}
沿用定理 \ref{thm:fold-gauge-anomaly-branch-radius-cumulant-growth} 的记号，记
\[
P_G(\theta)=\log\!\Bigl(\frac{\mu(e^\theta)}{2}\Bigr),
\qquad
\theta_\ast=\log u_\ast,
\qquad
R_\theta=|\theta_\ast|.
\]
则 \(P_G\) 的解析延拓在复平面上至少包含可数无穷多个奇点，并且包含两条平移晶格
\[
\theta_\ast+2\pi i\,\ZZ
\quad\text{与}\quad
\overline{\theta_\ast}+2\pi i\,\ZZ.
\]
因此 \(P_G\) 不是 \(D\)-有限（holonomic）函数。
\end{theorem}
\begin{proof}
由定理 \ref{thm:fold-gauge-anomaly-branch-radius-cumulant-growth}，\(\mu(u)\) 在 \(u=u_\ast\) 与 \(u=\overline{u_\ast}\) 处为平方根型分支点，故 \(P_G(\theta)=\log(\mu(e^\theta)/2)\) 在满足 \(e^\theta=u_\ast\) 或 \(e^\theta=\overline{u_\ast}\) 的点处不可解析。
\par
另一方面，由指数映射的 \(2\pi i\)-周期性，
\[
e^{\theta+2\pi i k}=e^\theta\qquad(k\in\ZZ),
\]
故
\[
e^\theta=u_\ast\iff \theta\in \log u_\ast+2\pi i\,\ZZ=\theta_\ast+2\pi i\,\ZZ,
\]
同理得到 \(\overline{\theta_\ast}+2\pi i\,\ZZ\)。
从而 \(P_G\) 的奇点集合至少包含上述两条晶格，因而为可数无穷。
\par
若 \(P_G\) 为 \(D\)-有限，则其满足某个首项系数为多项式的线性常微分方程；
经典理论表明其在复平面上的有限奇点只能落在首项系数多项式的零点集合中，因而有限奇点至多有限多个。
与上面的无穷奇点晶格矛盾，故 \(P_G\) 非 \(D\)-有限。证毕。
\end{proof}

\begin{corollary}[累积量密度序列的非 \texorpdfstring{$P$}{P}-递归性]\label{cor:fold-gauge-anomaly-cumulants-not-precursive}
记 \(\kappa_n:=P_G^{(n)}(0)\)（\(n\ge 0\)）。
则序列 \((\kappa_n)_{n\ge 0}\) 不是 \(P\)-递归的：不存在多项式 \(q_0,\dots,q_s\in\CC[n]\)（不全为零）使对充分大 \(n\) 有
\[
q_s(n)\kappa_{n+s}+\cdots+q_0(n)\kappa_n=0.
\]
\end{corollary}
\begin{proof}
若 \((\kappa_n)\) 为 \(P\)-递归，则其指数型生成函数
\[
\sum_{n\ge 0}\kappa_n\frac{\theta^n}{n!}=P_G(\theta)
\]
为 \(D\)-有限函数（\(P\)-递归与指数型生成函数 \(D\)-有限性之间的标准等价）。
这与定理 \ref{thm:fold-gauge-anomaly-singularity-lattice-nonholonomic} 得到的 \(P_G\) 非 \(D\)-有限矛盾。证毕。
\end{proof}

\begin{proposition}[根单位采样的最小收敛模数阈值]\label{prop:fold-gauge-anomaly-root-unit-sampling-threshold}
沿用定理 \ref{thm:fold-gauge-anomaly-branch-radius-cumulant-growth} 的解析半径 \(R_\theta\)。
对每个整数 \(p\ge 1\) 定义根单位采样点
\[
\theta_p:=\frac{2\pi i}{p}.
\]
则 \(P_G\) 在 \(\theta=0\) 的 Taylor 展开在 \(\theta=\theta_p\) 处收敛当且仅当
\[
|\theta_p|<R_\theta
\quad\Longleftrightarrow\quad
p>\frac{2\pi}{R_\theta}.
\]
特别地，取定理 \ref{thm:fold-gauge-anomaly-branch-radius-cumulant-growth} 中的数值
\(
R_\theta=0.7439369247693026
\)，则满足收敛的最小整数为 \(p_{\min}=9\)。
\end{proposition}
\begin{proof}
由复分析基本定理，\(P_G\) 在 \(\theta=0\) 的 Taylor 级数收敛圆盘半径恰等于最近奇点到原点的距离，亦即 \(R_\theta\)。
而 \(|\theta_p|=2\pi/p\)，故收敛当且仅当 \(2\pi/p<R_\theta\)，等价于 \(p>2\pi/R_\theta\)。
代入 \(R_\theta\) 的数值即可得 \(2\pi/R_\theta\approx 8.44\)，从而最小整数为 \(9\)。证毕。
\end{proof}

\endinput
